% Auto-generated from raw.txt
% Usage:
%   \vocabentry{word}{/ipa/ (pos)}{definition}{example}{note}

\vocabentry{Broadcasting}
{/ˈbrɔːdkæstɪŋ/ (n)}
{The distribution of audio or video content to a dispersed audience via any electronic mass communications medium.}
{The BBC is one of the oldest and most respected organizations in world broadcasting.}
{Đây là một danh từ dùng để chỉ ngành phát thanh hoặc truyền hình.}

\vocabentry{Mainstream media}
{/ˌmeɪnstriːm ˈmiːdiə/ (n)}
{Media disseminated via the largest distribution channels, which are therefore representative of what the majority of people consume.}
{Many people are turning to social media because they no longer trust the mainstream media.}
{Đây là một cụm danh từ dùng để chỉ truyền thông chính thống.}

\vocabentry{Endorsement}
{/ɪnˈdɔːrsmənt/ (n)}
{An act of giving one's public approval or support to someone or something, often in advertising.}
{The athlete signed a multi-million dollar endorsement deal with a major sports brand.}
{Đây là một danh từ dùng để chỉ sự xác nhận hoặc quảng cáo của người nổi tiếng cho một sản phẩm.}

\vocabentry{Circulation}
{/ˌsɜːrkjəˈleɪʃn/ (n)}
{The number of copies of a newspaper or magazine sold of each issue.}
{Digital versions have led to a steady decline in the physical circulation of daily newspapers.}
{Đây là một danh từ dùng để chỉ tổng số bản in được phát hành của một tờ báo hoặc tạp chí.}

\vocabentry{Censorship}
{/ˈsensərʃɪp/ (n)}
{The suppression or prohibition of any parts of books, films, news, etc. that are considered obscene, politically unacceptable, or a threat to security.}
{Tight government censorship makes it difficult for journalists to report the truth.}
{Đây là một danh từ dùng để chỉ sự kiểm duyệt.}

\vocabentry{Persuasive}
{/pərˈsweɪsɪv/ (adj)}
{Good at persuading someone to do or believe something through reasoning or the use of temptation.}
{Advertisers use highly persuasive language and imagery to influence consumer behavior.}
{Đây là một tính từ dùng để chỉ tính thuyết phục, đặc biệt quan trọng trong quảng cáo.}

\vocabentry{Algorithm}
{/ˈælɡərɪðəm/ (n)}
{A set of rules to be followed in calculations or other problem-solving operations, especially by a computer.}
{The social media algorithm determines which advertisements are shown to you based on your interests.}
{Trong truyền thông, đây là thuật toán điều phối nội dung.}

\vocabentry{Target audience}
{/ˌtɑːrɡɪt ˈɔːdiəns/ (n)}
{A particular group at which a film, book, advertising campaign, etc., is aimed.}
{The brand's target audience is primarily young professionals living in urban areas.}
{Đây là một cụm danh từ dùng để chỉ khách hàng/khán giả mục tiêu.}

\vocabentry{Prime time}
{/ˈpraɪm taɪm/ (n)}
{The time at which a radio or television audience is expected to be at its highest.}
{Advertising slots during prime time are significantly more expensive than those in the morning.}
{Đây là một cụm danh từ dùng để chỉ "giờ vàng" phát sóng.}

\vocabentry{Publicity}
{/pʌbˈlɪsəti/ (n)}
{Notice or attention given to someone or something by the media.}
{The scandal gained a lot of negative publicity for the tech company.}
{Đây là một danh từ dùng để chỉ sự công khai, sự chú ý từ dư luận hoặc hoạt động quảng bá.}

\vocabentry{Sensationalism}
{/senˈseɪʃənəlɪzəm/ (n)}
{The use of exciting or shocking stories or language at the expense of accuracy, in order to provoke public interest or excitement.}
{Tabloid newspapers are often accused of using sensationalism to boost their sales.}
{Đây là một danh từ dùng để chỉ lối đưa tin giật gân.}

\vocabentry{Billboard}
{/ˈbɪlbɔːrd/ (n)}
{A large outdoor board for displaying advertisements.}
{The new highway is lined with colorful billboards promoting local businesses.}
{Đây là một danh từ dùng để chỉ biển quảng cáo ngoài trời khổ lớn.}

\vocabentry{Brand loyalty}
{/ˈbrænd ˌlɔɪəlti/ (n)}
{The tendency of some consumers to continue buying the same brand of goods rather than competing brands.}
{Exceptional customer service is key to building long-term brand loyalty.}
{Đây là một cụm danh từ dùng để chỉ sự trung thành với thương hiệu.}

\vocabentry{Clickbait}
{/ˈklɪkbeɪt/ (n)}
{Content whose main purpose is to attract attention and encourage visitors to click on a link to a particular web page.}
{I was annoyed when the news article turned out to be nothing but clickbait.}
{Đây là một danh từ dùng để chỉ mồi nhử nhấp chuột, tiêu đề câu view.}

\vocabentry{Consumerism}
{/kənˈsuːmərɪzəm/ (n)}
{The protection or promotion of the interests of consumers; also the preoccupation of society with the acquisition of consumer goods.}
{Critics argue that excessive advertising fuels a culture of mindless consumerism.}
{Đây là một danh từ dùng để chỉ chủ nghĩa tiêu thụ.}

\vocabentry{Disseminate}
{/dɪˈsemɪneɪt/ (v)}
{To spread (something, especially information) widely.}
{The internet has made it easier than ever to disseminate news to a global audience.}
{Đây là một động từ dùng để chỉ việc phổ biến hoặc lan truyền thông tin.}

\vocabentry{Eye-catching}
{/ˈaɪ ˌkætʃɪŋ/ (adj)}
{Immediately noticeable because it is particularly interesting, bright, or attractive.}
{The poster features an eye-catching design that stands out in a crowded street.}
{Đây là một tính từ dùng để chỉ sự bắt mắt, thu hút sự chú ý.}

\vocabentry{Jingle}
{/ˈdʒɪŋɡl/ (n)}
{A short slogan, verse, or tune that is designed to be easily remembered, especially as used in advertising.}
{That company's commercial has such a catchy jingle that I can't get it out of my head.}
{Đây là một danh từ dùng để chỉ đoạn nhạc quảng cáo ngắn.}

\vocabentry{Mass media}
{/ˌmæs ˈmiːdiə/ (n)}
{The various methods of communication that reach large numbers of people.}
{The mass media plays a crucial role in shaping public opinion during elections.}
{Đây là một cụm danh từ dùng để chỉ các phương tiện truyền thông đại chúng.}

\vocabentry{Plagiarism}
{/ˈpleɪdʒərɪzəm/ (n)}
{The practice of taking someone else's work or ideas and passing them off as one's own.}
{The journalist was fired after being accused of plagiarism in his latest column.}
{Đây là một danh từ dùng để chỉ sự đạo văn, vi phạm bản quyền truyền thông.}

\vocabentry{Commercial}
{/kəˈmɜːrʃl/ (n)}
{A television or radio advertisement.}
{During the Super Bowl, many companies pay millions for a 30-second commercial.}
{Đây là một danh từ dùng để chỉ đoạn quảng cáo trên truyền hình hoặc phát thanh.}

\vocabentry{Logo}
{/ˈloʊɡoʊ/ (n)}
{A symbol or other design adopted by an organization to identify its products.}
{A well-designed logo should be simple, memorable, and reflective of the brand's identity.}
{Đây là một danh từ dùng để chỉ biểu tượng thương hiệu.}

\vocabentry{Marketing}
{/ˈmɑːrkɪtɪŋ/ (n)}
{The action or business of promoting and selling products or services.}
{Digital marketing has become essential for businesses to reach younger consumers.}
{Đây là một danh từ dùng để chỉ ngành tiếp thị.}

\vocabentry{Propaganda}
{/ˌprɑːpəˈɡændə/ (n)}
{Information, especially of a biased or misleading nature, used to promote a political cause or point of view.}
{The documentary was criticized for being a piece of government propaganda.}
{Đây là một danh từ dùng để chỉ hoạt động tuyên truyền.}

\vocabentry{Slogan}
{/ˈsloʊɡən/ (n)}
{A short and striking or memorable phrase used in advertising.}
{The company's famous slogan, "Just Do It," is recognized worldwide.}
{Đây là một danh từ dùng để chỉ khẩu hiệu quảng cáo.}

\vocabentry{Subscription}
{/səbˈskrɪpʃn/ (n)}
{The action of making or agreeing to make an advance payment in order to receive or participate in something.}
{I have a monthly subscription to an online streaming service.}
{Đây là một danh từ dùng để chỉ việc đăng ký thuê bao báo chí hoặc dịch vụ trực tuyến.}

\vocabentry{Viral}
{/ˈvaɪrəl/ (adj)}
{Circulated rapidly and widely from one internet user to another.}
{A viral marketing campaign can be much more effective than traditional advertising.}
{Đây là một tính từ dùng để chỉ sự lan truyền nhanh chóng trên mạng, đặc biệt là nội dung được chia sẻ rộng rãi}

\vocabentry{Yellow journalism}
{/ˌjeloʊ ˈdʒɜːrnəlɪzəm/ (n)}
{Journalism that is based upon sensationalism and crude exaggeration.}
{The politician accused the newspaper of practicing yellow journalism to harm his reputation.}
{Đây là một cụm danh từ dùng để chỉ lối làm báo "lá cải", thiếu khách quan.}

\vocabentry{Banner ad}
{/ˈbænər æd/ (n)}
{An advertisement appearing on a web page in the form of a bar, column, or box.}
{Banner ads are often placed at the top or sides of a webpage to attract attention.}
{Đây là một cụm danh từ dùng để chỉ quảng cáo biểu ngữ trên trang web.}

\vocabentry{Coverage}
{/ˈkʌvərɪdʒ/ (n)}
{The treatment of an issue by the media.}
{The Olympic Games received extensive media coverage across the globe.}
{Đây là một danh từ dùng để chỉ việc đưa tin hoặc độ phủ sóng của tin tức.}

\vocabentry{Deadline}
{/ˈdedlaɪn/ (n)}
{The latest time or date by which something should be completed.}
{Journalists often have to work under pressure to meet tight deadlines.}
{Đây là một danh từ dùng để chỉ hạn chót nộp bài hoặc hoàn thành chiến dịch.}

\vocabentry{Editor}
{/ˈedɪtər/ (n)}
{A person who is in charge of and determines the final content of a newspaper, magazine, or multi-author book.}
{The editor decided to move the story to the front page of the newspaper.}
{Đây là một danh từ dùng để chỉ biên tập viên hoặc chủ bút.}

\vocabentry{Feature}
{/ˈfiːtʃər/ (n)}
{A newspaper or magazine article or a broadcast program devoted to a particular topic.}
{The Sunday magazine includes a special feature on sustainable living.}
{Đây là một danh từ dùng để chỉ bài viết chuyên sâu hoặc phóng sự đặc biệt.}

\vocabentry{Informed}
{/ɪnˈfɔːrmd/ (adj)}
{Having or showing knowledge of a particular subject or situation.}
{Mass media helps to create an informed citizenry that can participate in democracy.}
{Đây là một tính từ dùng để chỉ sự am hiểu, thường dùng cho "informed public" - công chúng có hiểu biết.}

\vocabentry{Journalism}
{/ˈdʒɜːrnəlɪzəm/ (n)}
{The activity or profession of writing for newspapers, magazines, or news websites or preparing news to be broadcast.}
{She decided to pursue a career in investigative journalism.}
{Đây là một danh từ dùng để chỉ ngành báo chí.}

\vocabentry{Mediate}
{/ˈmiːdieɪt/ (v)}
{To communicate or transmit (something) through a medium.}
{Our experience of the world is increasingly mediated through digital screens.}
{Trong truyền thông, đây là động từ chỉ việc truyền đạt thông tin qua các kênh trung gian.}

\vocabentry{Monopoly}
{/məˈnɑːpəli/ (n)}
{The exclusive possession or control of the supply of or trade in a commodity or service.}
{Critics are concerned about the media monopoly held by a few large corporations.}
{Trong truyền thông, đây là sự độc quyền truyền thông.}

\vocabentry{Opinion}
{/əˈpɪnjən/ (n)}
{A view or judgment formed about something, not necessarily based on fact or knowledge.}
{Columnists are paid to express their personal opinions on current events.}
{Đây là một danh từ dùng để chỉ quan điểm, thường thấy trong các mục "Opinion piece".}

\vocabentry{Press release}
{/ˈpres rɪˌliːs/ (n)}
{An official statement issued to newspapers giving information on a particular matter.}
{The company issued a press release to announce the launch of its new product.}
{Đây là một cụm danh từ dùng để chỉ thông cáo báo chí.}

\vocabentry{Ratings}
{/ˈreɪtɪŋz/ (n)}
{A measurement of the size of the audience for a television or radio program.}
{The show was canceled due to low ratings and poor critical reviews.}
{Đây là một danh từ dùng để chỉ tỉ lệ người xem/nghe đài.}

\vocabentry{Script}
{/skrɪpt/ (n)}
{The written text of a play, film, or broadcast.}
{The scriptwriters worked for months to perfect the dialogue for the new movie.}
{Đây là một danh từ dùng để chỉ kịch bản.}

\vocabentry{Source}
{/sɔːrs/ (v)}
{Obtain from a particular source.}
{The reporter struggled to source reliable information about the secret government project.}
{Trong báo chí, đây là việc tìm kiếm nguồn tin.}

\vocabentry{Sponsorship}
{/ˈspɑːnsərʃɪp/ (n)}
{The position of being a sponsor.}
{The local soccer team relies on corporate sponsorship to cover their travel costs.}
{Đây là một danh từ dùng để chỉ sự tài trợ.}

\vocabentry{Tabloid}
{/ˈtæblɔɪd/ (n)}
{A newspaper having pages half the size of those of a standard newspaper, typically popular in style and dominated by headlines, photographs, and sensational stories.}
{Tabloids are known for focusing on celebrity gossip and scandals.}
{Đây là một danh từ dùng để chỉ báo khổ nhỏ/báo lá cải.}

\vocabentry{Telecommunication}
{/ˌtelikəˌmjuːnɪˈkeɪʃn/ (n)}
{Communication over a distance by cable, telegraph, telephone, or broadcasting.}
{Advances in telecommunication have transformed the way we receive news.}
{Đây là một danh từ dùng để chỉ viễn thông.}

\vocabentry{Underestimate}
{/ˌʌndərˈestɪmeɪt/ (v)}
{Estimate (something) to be smaller or less important than it actually is.}
{Politicians often underestimate the power of social media to influence voters.}
{Trong truyền thông, thường nói về việc đánh giá thấp sức mạnh của dư luận.}

\vocabentry{Viewer}
{/ˈvjuːər/ (n)}
{A person who watches television or a movie.}
{The show was a hit with viewers and critics alike.}
{Đây là một danh từ dùng để chỉ khán giả, người xem truyền hình hoặc phim}

\vocabentry{Webcam}
{/ˈwebkæm/ (n)}
{A video camera that inputs to a computer connected to the internet, so that its images can be viewed by internet users.}
{Journalists use webcams to conduct remote interviews with experts around the world.}
{Đây là một danh từ liên quan đến truyền thông trực tuyến.}

\vocabentry{Whistleblower}
{/ˈwɪslbloʊər/ (n)}
{A person who informs on a person or organization engaged in an illicit activity.}
{The whistleblower leaked confidential documents to the press to expose the company's corruption.}
{Đây là một danh từ dùng để chỉ người tố giác, nguồn tin quan trọng cho báo chí điều tra.}

\vocabentry{Wireless}
{/ˈwaɪərləs/ (adj)}
{Using radio, microwaves, etc., as opposed to wires or cables to transmit signals.}
{Wireless headphones are very convenient for listening to podcasts on the go.}
{Đây là một tính từ dùng để chỉ công nghệ không dây, truyền tải thông tin không cần cáp vật lý}

\vocabentry{Aggregator}
{/ˈæɡrɪɡeɪtər/ (n)}
{A website or program that collects related items of content and displays them.}
{A news aggregator allows you to see headlines from many different sources in one place.}
{Đây là một danh từ dùng để chỉ trình tổng hợp tin tức.}

\vocabentry{Bias}
{/ˈbaɪəs/ (n)}
{Prejudice in favor of or against one thing, person, or group compared with another, usually in a way considered to be unfair.}
{Many viewers complain about the political bias in certain news networks.}
{Đây là một danh từ dùng để chỉ sự thiên kiến, mất khách quan.}

\vocabentry{Brief}
{/briːf/ (n)}
{A set of instructions given to a person about a job or task.}
{The creative team is working on a new campaign based on the client's brief.}
{Trong quảng cáo, đây là bản tóm tắt yêu cầu sáng tạo từ khách hàng.}

\vocabentry{Broadsheet}
{/ˈbrɔːdʃiːt/ (n)}
{A newspaper with a large format, regarded as more serious and less sensational than tabloids.}
{Broadsheets tend to focus on in-depth analysis of politics and business.}
{Đây là một danh từ dùng để chỉ báo khổ lớn, chính thống.}

\vocabentry{Caption}
{/ˈkæpʃn/ (n)}
{A title or brief explanation accompanying an illustration, cartoon, or poster.}
{An effective caption should provide context for the image and engage the reader.}
{Đây là một danh từ dùng để chỉ lời chú thích cho ảnh hoặc video.}

\vocabentry{Categorize}
{/ˈkætəɡəraɪz/ (v)}
{Place in a particular class or group.}
{News websites categorize stories into sections like World, Business, and Sports.}
{Trong truyền thông, đây là việc phân loại tin tức/quảng cáo.}

\vocabentry{Celebrity}
{/səˈlebrəti/ (n)}
{A famous person.}
{Advertisers often use celebrity endorsements to increase brand awareness.}
{Đây là một danh từ dùng để chỉ người nổi tiếng.}

\vocabentry{Column}
{/ˈkɑːləm/ (n)}
{A vertical division of a page or text; also a regular section of a newspaper or magazine.}
{She writes a weekly column on fashion for the city's main newspaper.}
{Đây là một danh từ dùng để chỉ cột báo hoặc chuyên mục.}

\vocabentry{Commentator}
{/ˈkɑːmənˌteɪtər/ (n)}
{A person who comments on events or on a text.}
{Political commentators are analyzing the impact of the new law on the economy.}
{Đây là một danh từ dùng để chỉ bình luận viên.}

\vocabentry{Compact}
{/kəmˈpækt/ (adj)}
{Closely and neatly packed together.}
{This compact camera is perfect for mobile journalists on the move.}
{Thường dùng cho các thiết bị truyền thông nhỏ gọn.}

\vocabentry{Conflict of interest}
{/ˌkɑːnflɪkt əv ˈɪntrest/ (n)}
{A situation in which a person is in a position to derive personal benefit from actions or decisions made in their official capacity.}
{The journalist was criticized for a conflict of interest when he wrote about his brother's company.}
{Đây là một cụm danh từ dùng để chỉ sự xung đột lợi ích.}

\vocabentry{Conglomerate}
{/kənˈɡlɑːmərət/ (n)}
{A large corporation formed by the merging of separate and diverse firms.}
{The media conglomerate owns several television stations, newspapers, and film studios.}
{Trong truyền thông, đây là các tập đoàn truyền thông khổng lồ.}

\vocabentry{Credibility}
{/ˌkredəˈbɪləti/ (n)}
{The quality of being trusted and believed in.}
{A news organization's credibility depends on its commitment to accuracy and fairness.}
{Đây là một danh từ dùng để chỉ sự uy tín, độ tin cậy.}

\vocabentry{Critic}
{/ˈkrɪtɪk/ (n)}
{A person who expresses an unfavorable opinion of something.}
{Film critics gave the new blockbuster very positive reviews.}
{Đây là một danh từ dùng để chỉ nhà phê bình.}

\vocabentry{Data protection}
{/ˈdeɪtə prəˌtekʃn/ (n)}
{Legal control over access to and use of data stored in computers.}
{New regulations have been introduced to strengthen data protection for internet users.}
{Đây là một cụm danh từ dùng để chỉ bảo mật dữ liệu.}

\vocabentry{Debate}
{/dɪˈbeɪt/ (n)}
{A formal discussion on a particular topic in a public meeting or legislative assembly.}
{The televised debate between the two candidates attracted millions of viewers.}
{Đây là một danh từ dùng để chỉ cuộc tranh luận.}

\vocabentry{Demographic}
{/ˌdeməˈɡræfɪk/ (n)}
{A particular sector of a population.}
{The advertisement is aimed at the 18-34 age demographic.}
{Trong quảng cáo, đây là phân khúc nhân khẩu học.}

\vocabentry{Diverse}
{/daɪˈvɜːrs/ (adj)}
{Showing a great deal of variety; very different.}
{It is important for the media to reflect the diverse voices and perspectives of society.}
{Đây là một tính từ dùng để chỉ sự đa dạng trong nội dung truyền thông.}

\vocabentry{Documentary}
{/ˌdɑːkjuˈmentri/ (n)}
{A movie or a television or radio program that provides a factual record or report.}
{The documentary explores the impact of plastic pollution on marine life.}
{Đây là một danh từ dùng để chỉ phim tài liệu.}

\vocabentry{Edit}
{/ˈedɪt/ (v)}
{Prepare (written material, audio, or video) for publication by correcting, condensing, or otherwise modifying it.}
{It took several hours to edit the raw footage into a three-minute news report.}
{Đây là một động từ dùng để chỉ việc biên tập.}

\vocabentry{Exclusive}
{/ɪkˈskluːsɪv/ (n)}
{A story or interview that is published by only one newspaper or broadcast by only one television station.}
{The magazine paid a large sum for an exclusive interview with the royal family.}
{Đây là một danh từ dùng để chỉ tin độc quyền.}

\vocabentry{Expose}
{/ɪkˈspoʊz/ (v)}
{Make (something) visible by uncovering it; in journalism, to reveal a scandal.}
{The investigation aimed to expose the corrupt practices of the local government.}
{Đây là một động từ dùng để chỉ việc phơi bày hoặc vạch trần.}

\vocabentry{Fake news}
{/ˌfeɪk ˈnuːz/ (n)}
{False stories that appear to be news spread on the internet or using other media.}
{Fake news can spread quickly on social media and cause widespread confusion.}
{Đây là một cụm danh từ dùng để chỉ tin giả.}

\vocabentry{Feedback}
{/ˈfiːdbæk/ (n)}
{Information about reactions to a product, a person's performance of a task, etc.}
{The television station welcomes feedback from its viewers on its programming.}
{Đây là một danh từ dùng để chỉ phản hồi từ khán giả/khách hàng.}

\vocabentry{Flyer}
{/ˈflaɪər/ (n)}
{A small handbill advertising an event or product.}
{I was handed a flyer on the street advertising a new restaurant opening nearby.}
{Đây là một danh từ dùng để chỉ tờ rơi quảng cáo.}

\vocabentry{Follower}
{/ˈfɑːloʊər/ (n)}
{A person who subscribes to a person or group on social media.}
{The influencer has over a million followers on her travel blog.}
{Đây là một danh từ dùng để chỉ người theo dõi trên mạng xã hội.}

\vocabentry{Front page}
{/ˌfrʌnt ˈpeɪdʒ/ (n)}
{The first page of a newspaper, containing the most important stories.}
{The earthquake was the lead story on the front page of every national newspaper.}
{Đây là một cụm danh từ dùng để chỉ trang nhất.}

\vocabentry{Gossip}
{/ˈɡɑːsɪp/ (n)}
{Casual or unconstrained conversation or reports about other people, typically involving details that are not confirmed as being true.}
{Tabloid newspapers thrive on celebrity gossip and rumors.}
{Đây là một danh từ dùng để chỉ tin đồn, chuyện phiếm.}

\vocabentry{Headline}
{/ˈhedlaɪn/ (n)}
{A heading at the top of an article or page in a newspaper or magazine.}
{The headline was bold and attention-grabbing to encourage people to read the story.}
{Đây là một danh từ dùng để chỉ tiêu đề bài báo.}

\vocabentry{Hype}
{/haɪp/ (n)}
{Extravagant or intensive publicity or promotion.}
{There was a lot of hype surrounding the release of the new smartphone.}
{Đây là một danh từ dùng để chỉ sự cường điệu hoặc quảng cáo thổi phồng.}

\vocabentry{Icon}
{/ˈaɪkɑːn/ (n)}
{A person or thing regarded as a representative symbol or as worthy of veneration.}
{Marilyn Monroe remains a global fashion and beauty icon.}
{Trong truyền thông, đây là biểu tượng văn hóa.}

\vocabentry{Illustrate}
{/ˈɪləstreɪt/ (v)}
{Provide (a book, newspaper, etc.) with pictures.}
{The article is illustrated with several high-quality photographs and charts.}
{Đây là một động từ dùng để chỉ việc minh họa.}

\vocabentry{Impression}
{/ɪmˈpreʃn/ (n)}
{An instance of a pop-up or other ad being displayed on a user's screen.}
{The marketing campaign generated over a million impressions in its first week.}
{Trong quảng cáo trực tuyến, đây là lượt hiển thị.}

\vocabentry{Influencer}
{/ˈɪnfluənsər/ (n)}
{A person with the ability to influence potential buyers of a product or service by promoting or recommending the items on social media.}
{Many brands are collaborating with influencers to reach younger audiences.}
{Đây là một danh từ dùng để chỉ người có sức ảnh hưởng.}

\vocabentry{Information age}
{/ˌɪnfərˈmeɪʃn eɪdʒ/ (n)}
{The period of time during which large amounts of information became widely available to many people.}
{We are living in the information age where data is more valuable than oil.}
{Đây là một cụm danh từ dùng để chỉ kỷ nguyên thông tin.}

\vocabentry{Infomercial}
{/ˌɪnfoʊˈmɜːrʃl/ (n)}
{A television program that is an extended advertisement.}
{The late-night infomercial promoted a new range of fitness equipment.}
{Đây là một danh từ dùng để chỉ chương trình quảng cáo mang tính cung cấp thông tin.}

\vocabentry{In-depth}
{/ˌɪn ˈdepθ/ (adj)}
{Comprehensive and thorough.}
{The magazine provides in-depth analysis of global economic trends.}
{Đây là một tính từ dùng để chỉ tính sâu sắc, kỹ lưỡng của một bài báo cáo.}

\vocabentry{Investigative}
{/ɪnˈvestɪɡeɪtɪv/ (adj)}
{Designed to find out the truth about something, especially a crime or a scandal. (Đây là một tính từ dùng để chỉ tính điều tra (trong báo chí điều tra).).}
{Investigative reporters spent months uncovering the financial scandal.}
{}

\vocabentry{Journalist}
{/ˈdʒɜːrnəlɪst/ (n)}
{A person who writes for newspapers, magazines, or news websites or prepares news to be broadcast.}
{A journalist's job is to report the facts accurately and objectively.}
{Đây là một danh từ dùng để chỉ nhà báo.}

\vocabentry{Knowledge}
{/ˈnɑːlɪdʒ/ (n)}
{Facts, information, and skills acquired by a person through experience or education.}
{Access to reliable information is essential for gaining knowledge about the world.}
{Truyền thông là công cụ truyền tải kiến thức.}

\vocabentry{Leak}
{/liːk/ (v)}
{Intentionally disclose (secret information).}
{Someone in the government leaked the confidential report to the press.}
{Đây là một động từ dùng để chỉ việc rò rỉ thông tin bí mật.}

\vocabentry{Libel}
{/ˈlaɪbl/ (n)}
{A published false statement that is damaging to a person's reputation; a written defamation.}
{The actor sued the newspaper for libel after they published a false story about him.}
{Đây là một danh từ dùng để chỉ tội phỉ báng bằng văn bản.}

\vocabentry{Link}
{/lɪŋk/ (n)}
{A reference to data that the reader can directly follow.}
{You can click on the link to read the full article on the news website.}
{Đây là một danh từ dùng để chỉ đường liên kết trực tuyến.}

\vocabentry{Live}
{/laɪv/ (adj)}
{(Of a performance or broadcast) seen or heard at the same time as it happens. (Đây là một tính từ dùng để chỉ sự trực tiếp (phát sóng trực tiếp).).}
{The television station provided live coverage of the royal wedding.}
{}

\vocabentry{Log}
{/lɔːɡ/ (v)}
{Make a systematic record of.}
{Websites often log user activity to personalize their experience.}
{Trong truyền thông, đây là việc ghi nhật ký phát sóng hoặc dữ liệu người dùng.}

\vocabentry{Mainstream}
{/ˈmeɪnstriːm/ (adj)}
{Representing widespread current thought, influence, or activity.}
{Mainstream news outlets often avoid reporting on niche or controversial topics.}
{Đây là một tính từ dùng để chỉ tính chủ lưu, dòng chính của truyền thông hoặc văn hóa đại chúng}

\vocabentry{Manipulation}
{/məˌnɪpjuˈleɪʃn/ (n)}
{The action of manipulating someone or something in a clever or unscrupulous way.}
{Critics are concerned about the manipulation of public opinion through social media bots.}
{Đây là một danh từ dùng để chỉ sự thao túng, thường dùng trong "media manipulation".}

\vocabentry{Media literacy}
{/ˌmiːdiə ˈlɪtərəsi/ (n)}
{The ability to access, analyze, evaluate, and create media in a variety of forms.}
{Schools are increasingly teaching media literacy to help students spot fake news.}
{Đây là một cụm danh từ dùng để chỉ kiến thức truyền thông, khả năng nhận diện tin tức.}

\vocabentry{Message}
{/ˈmesɪdʒ/ (n)}
{A significant political or social points that an author or artist tries to convey.}
{The main message of the campaign is that everyone has the power to protect the planet.}
{Trong quảng cáo, đây là thông điệp cốt lõi.}

\vocabentry{Misinformation}
{/ˌmɪsɪnfərˈmeɪʃn/ (n)}
{False or inaccurate information that is spread, regardless of intent to deceive.}
{The spread of misinformation about vaccines has become a major public health concern.}
{Đây là một danh từ dùng để chỉ thông tin sai lệch.}

\vocabentry{Network}
{/ˈnetwɜːrk/ (n)}
{A group of interconnected people or things; a group of broadcasting stations.}
{The television network has affiliates in almost every major city.}
{Đây là một danh từ dùng để chỉ mạng lưới truyền hình/phát thanh.}

\vocabentry{Newsfeed}
{/ˈnuːzfiːd/ (n)}
{A service that provides a constantly updated list of news stories.}
{I check my social media newsfeed every morning for the latest updates.}
{Đây là một danh từ dùng để chỉ bảng tin cập nhật trên trang web/mạng xã hội.}

\vocabentry{Newsflash}
{/ˈnuːzflæʃ/ (n)}
{A short news report, which is broadcast during another program when something very important has happened.}
{We interrupt this program for a special newsflash about the earthquake.}
{Đây là một danh từ dùng để chỉ tin vắn, tin khẩn cấp.}

\vocabentry{Newspaper}
{/ˈnuːzpeɪpər/ (n)}
{A printed publication (usually issued daily or weekly) consisting of folded unstapled sheets and containing news, feature articles, advertisements, and correspondence.}
{Reading a newspaper is a great way to stay informed about local and international news.}
{Đây là một danh từ dùng để chỉ tờ báo.}

\vocabentry{Objective}
{/əbˈdʒektɪv/ (adj)}
{(Of a person or their judgment) not influenced by personal feelings or opinions in considering and representing facts.}
{A journalist must remain objective and present all sides of a story.}
{Đây là một tính từ dùng để chỉ tính khách quan trong báo chí.}

\vocabentry{Observe}
{/əbˈzɜːrv/ (v)}
{Watch (someone or something) carefully and attentively.}
{To write a good feature, you must carefully observe your subject and their environment.}
{Phóng viên cần quan sát kỹ thực tế.}

\vocabentry{On-air}
{/ˌɑːn ˈer/ (adj/adv)}
{Being broadcast on radio or television.}
{The presenter didn't realize she was still on-air when she made the joke.}
{Đây là một cụm từ dùng để chỉ trạng thái đang lên sóng.}

\vocabentry{Outline}
{/ˈaʊtlaɪn/ (n)}
{A general description or plan giving the essential features of something but not the detail.}
{The journalist prepared a brief outline of the story before he started writing.}
{Trong viết lách truyền thông, đây là dàn ý bài viết.}

\vocabentry{Outreach}
{/ˈaʊtriːtʃ/ (n)}
{The extent or length of reaching out; also an organization's involvement with the community.}
{The museum's educational outreach program aims to engage students from low-income areas.}
{Trong PR, đây là hoạt động tiếp cận cộng đồng.}

\vocabentry{Overload}
{/ˈoʊvərloʊd/ (n)}
{An excessive amount of something.}
{Information overload in the digital age can lead to stress and decision fatigue.}
{Thường dùng: "information overload" - quá tải thông tin.}

\vocabentry{Paper}
{/ˈpeɪpər/ (n)}
{A newspaper.}
{Have you seen today's paper? There's a big story about the new bridge.}
{Đây là một danh từ không trang trọng dùng để chỉ tờ báo.}

\vocabentry{Parody}
{/ˈpærədi/ (n)}
{An imitation of the style of a particular writer, artist, or genre with deliberate exaggeration for comic effect.}
{The comedian performed a hilarious parody of a typical news broadcast.}
{Đây là một danh từ dùng để chỉ tác phẩm nhại lại, châm biếm.}

\vocabentry{Periodical}
{/ˌpɪriˈɑːdɪkl/ (n)}
{A magazine or newspaper published at regular intervals.}
{The library has a wide selection of academic periodicals and journals.}
{Đây là một danh từ dùng để chỉ ấn phẩm định kỳ.}

\vocabentry{Persuasion}
{/pərˈsweɪʒn/ (n)}
{The action or fact of persuading someone or of being persuaded to do or believe something.}
{Subtle persuasion is often more effective in advertising than aggressive sales tactics.}
{Quảng cáo là nghệ thuật của sự thuyết phục.}

\vocabentry{Phenomenon}
{/fəˈnɑːmɪnən/ (n)}
{A fact or situation that is observed to exist or happen, especially one whose cause or explanation is in question.}
{The rapid rise of TikTok is a fascinating cultural phenomenon.}
{Sự nổi tiếng trên mạng là một hiện tượng xã hội.}

\vocabentry{Podcast}
{/ˈpɑːdkæst/ (n)}
{A digital audio file made available on the internet for downloading to a computer or mobile device.}
{I listen to several podcasts about history and science during my commute.}
{Đây là một danh từ dùng để chỉ tệp âm thanh kỹ thuật số.}

\vocabentry{Post}
{/poʊst/ (n)}
{An item of content published on social media or a blog.}
{Her latest post about her vacation received hundreds of likes and comments.}
{Đây là một danh từ dùng để chỉ bài đăng.}

\vocabentry{Presenter}
{/prɪˈzentər/ (n)}
{A person who introduces and appears in a television or radio program.}
{He has been the main news presenter for the national station for over ten years.}
{Đây là một danh từ dùng để chỉ người dẫn chương trình.}

\vocabentry{Press}
{/pres/ (n)}
{Newspapers and magazines collectively.}
{The freedom of the press is a fundamental right in many democratic countries.}
{Đây là một danh từ dùng để chỉ giới báo chí.}

\vocabentry{Print}
{/prɪnt/ (n)}
{The production of books, newspapers, or other printed material.}
{Despite the rise of digital media, many people still prefer reading in print.}
{Đây là một danh từ dùng để chỉ ngành in ấn/báo in.}

\vocabentry{Privacy policy}
{/ˈpraɪvəsi ˌpɑːləsi/ (n)}
{A statement or legal document that discloses some or all of the ways a party gathers, uses, discloses, and manages a customer's or client's data.}
{You should always read a website's privacy policy before sharing your personal information.}
{Đây là một cụm danh từ dùng để chỉ chính sách quyền riêng tư.}

\vocabentry{Product placement}
{/ˈprɑːdʌkt ˌpleɪsmənt/ (n)}
{A practice in which manufacturers of goods or providers of a service gain exposure for their products by paying for them to be featured in movies and television programs.}
{The James Bond movies are famous for their subtle and not-so-subtle product placement.}
{Đây là một cụm danh từ dùng để chỉ quảng cáo chèn sản phẩm trong phim/chương trình.}

\vocabentry{Profile}
{/ˈproʊfaɪl/ (n)}
{A short article giving a description of a person or organization.}
{The magazine published a fascinating profile of the reclusive billionaire.}
{Đây là một danh từ dùng để chỉ bài giới thiệu nhân vật hoặc hồ sơ cá nhân.}

\vocabentry{Promote}
{/prəˈmoʊt/ (v)}
{Give publicity to (a product, organization, or venture) so as to increase sales or public awareness.}
{The pop star is traveling around the world to promote her new album.}
{Đây là một động từ dùng để chỉ việc quảng bá.}

\vocabentry{Publication}
{/ˌpʌblɪˈkeɪʃn/ (n)}
{The preparation and issuing of a book, journal, piece of music, or other work for public sale.}
{The university has its own publication for student research and creative writing.}
{Đây là một danh từ dùng để chỉ sự xuất bản hoặc ấn phẩm.}

\vocabentry{Public relations}
{/ˌpʌblɪk rɪˈleɪʃnz/ (n)}
{The professional maintenance of a favorable public image by a company or other organization or a famous person.}
{A good public relations strategy can help a company rebuild its reputation after a scandal.}
{Đây là một cụm danh từ dùng để chỉ ngành quan hệ công chúng - PR.}

\vocabentry{Publish}
{/ˈpʌblɪʃ/ (v)}
{Prepare and issue for public sale.}
{The newspaper decided not to publish the controversial photos.}
{Đây là một động từ dùng để chỉ việc xuất bản.}

\vocabentry{Reach}
{/riːtʃ/ (n)}
{The total number of different people or households exposed, at least once, to a medium during a given period.}
{Social media allows small businesses to have a much larger reach than traditional advertising.}
{Trong quảng cáo, đây là lượt tiếp cận.}

\vocabentry{Reaction}
{/riˈækʃn/ (n)}
{An action performed or a feeling experienced in response to a situation or event.}
{The public's reaction to the new advertisement was overwhelmingly positive.}
{Phản ứng của công chúng là thước đo truyền thông.}

\vocabentry{Reader}
{/ˈriːdər/ (n)}
{A person who reads.}
{Magazines often conduct surveys to find out more about their regular readers.}
{Đây là một danh từ dùng để chỉ độc giả.}

\vocabentry{Reliable}
{/rɪˈlaɪəbl/ (adj)}
{Consistently good in quality or performance; able to be trusted.}
{It is important to find reliable news sources that report facts accurately.}
{Đây là một tính từ dùng để chỉ sự đáng tin cậy của thông tin.}

\vocabentry{Reporter}
{/rɪˈpɔːrtər/ (n)}
{A person who reports news or conducts interviews for newspapers or broadcasts.}
{The reporter was on the scene within minutes of the accident.}
{Đây là một danh từ dùng để chỉ phóng viên.}

\vocabentry{Retract}
{/rɪˈtrækt/ (v)}
{Withdraw (a statement or accusation) as untrue or unjustified.}
{The newspaper was forced to retract the story and issue a formal apology.}
{Đây là một động từ dùng để chỉ việc rút lại lời tuyên bố hoặc tin đã đăng.}

\vocabentry{Review}
{/rɪˈvjuː/ (n)}
{A formal assessment or examination of something with the possibility or intention of instituting change if necessary.}
{I always read the movie reviews before I decide which film to see.}
{Đây là một danh từ dùng để chỉ bài đánh giá, phê bình.}

\vocabentry{Scoop}
{/skuːp/ (n)}
{A piece of news published by a newspaper or broadcast by a television station before its rivals.}
{The young journalist got a major scoop when she interviewed the reclusive actor.}
{Đây là một danh từ dùng để chỉ tin sốt dẻo, tin nhanh hơn đối thủ.}

\vocabentry{Search engine}
{/ˈsɜːrtʃ ˌendʒɪn/ (n)}
{A program that searches for and identifies items in a database that correspond to keywords.}
{Most internet users use a search engine to find information and products online.}
{Công cụ quan trọng trong tiếp thị số.}

\vocabentry{Section}
{/ˈsekʃn/ (n)}
{Any of the more or less distinct parts into which something is or may be divided.}
{I usually skip the business section and go straight to the sports news.}
{Đây là một danh từ dùng để chỉ các chuyên mục của tờ báo.}

\vocabentry{Sequel}
{/ˈsiːkwəl/ (n)}
{A published, broadcast, or recorded work that continues the story or develops the theme of an earlier one.}
{Fans are eagerly awaiting the sequel to the popular fantasy movie.}
{Đây là một danh từ dùng để chỉ phần tiếp theo.}

\vocabentry{Series}
{/ˈsɪriːz/ (n)}
{A set of related television or radio programs.}
{The documentary series explores the history of ancient civilizations.}
{Đây là một danh từ dùng để chỉ phim bộ hoặc chuỗi chương trình.}

\vocabentry{Short}
{/ʃɔːrt/ (n)}
{A short film.}
{The film festival features a competition for the best animated short.}
{Đây là một danh từ dùng để chỉ phim ngắn.}

\vocabentry{Small screen}
{/ˌsmɔːl ˈskriːn/ (n)}
{Television, especially in contrast to cinema.}
{Many famous movie stars are now appearing in high-quality dramas on the small screen.}
{Đây là một cụm danh từ dùng để chỉ màn ảnh nhỏ/truyền hình.}

\vocabentry{Soap opera}
{/ˈsoʊp ˌɑːprə/ (n)}
{A television or radio drama series dealing typically with daily events in the lives of the same group of characters.}
{My grandmother has been watching the same soap opera for over twenty years.}
{Đây là một danh từ dùng để chỉ phim truyền hình dài tập, kịch xà phòng.}

\vocabentry{Social network}
{/ˌsoʊʃl ˈnetwɜːrk/ (n)}
{A dedicated website or other application which enables users to communicate with each other.}
{Social networks have transformed the way we communicate and share information.}
{Đây là một cụm danh từ dùng để chỉ mạng xã hội.}

\vocabentry{Space}
{/speɪs/ (n)}
{An area or portion that is available for use.}
{Companies pay a lot of money for advertising space on the busiest websites.}
{Trong quảng cáo, đây là vị trí quảng cáo trên báo/web.}

\vocabentry{Spearhead}
{/ˈspɪrhed/ (v)}
{Lead (an attack or movement).}
{The marketing director will spearhead the launch of the new global campaign.}
{Đây là một động từ dùng để chỉ việc dẫn đầu một chiến dịch truyền thông.}

\vocabentry{Sponsor}
{/ˈspɑːnsər/ (n)}
{A person or organization that provides funds for a project or activity carried out by another.}
{The local art gallery is looking for sponsors for its next major exhibition.}
{Đây là một danh từ dùng để chỉ nhà tài trợ.}

\vocabentry{Sportscast}
{/ˈspɔːrtskæst/ (n)}
{A television or radio broadcast of sporting events.}
{I listen to the sportscast every morning to hear the latest scores and news.}
{Đây là một danh từ dùng để chỉ chương trình tin tức thể thao.}

\vocabentry{Spot}
{/spɑːt/ (n)}
{A short advertisement in a television or radio program.}
{The company bought a prime-time advertising spot during the movie.}
{Đây là một danh từ dùng để chỉ thời lượng quảng cáo ngắn.}

\vocabentry{Spread}
{/spred/ (v)}
{Extend over a large or increasing area.}
{Social media allows news to spread around the world in a matter of seconds.}
{Đây là một động từ dùng để chỉ việc lan truyền tin tức.}

\vocabentry{Stance}
{/stæns/ (n)}
{The attitude of a person or organization toward something.}
{The newspaper is known for its liberal stance on social and environmental issues.}
{Lập trường truyền thông của một tờ báo.}

\vocabentry{Station}
{/ˈsteɪʃn/ (n)}
{A company that broadcasts radio or television programs.}
{My favorite radio station plays a mix of classic rock and new music.}
{Đây là một danh từ dùng để chỉ đài truyền hình/phát thanh.}

\vocabentry{Stereotype}
{/ˈsteriətaɪp/ (n)}
{A widely held but fixed and oversimplified image or idea of a particular type of person or thing.}
{Critics argue that advertising often reinforces negative stereotypes about women.}
{Trong truyền thông, đây là khuôn mẫu định kiến.}

\vocabentry{Subliminal}
{/ˌsʌbˈlɪmɪnl/ (adj)}
{(Of a stimulus or mental process) below the threshold of sensation or consciousness.}
{Subliminal advertising is banned in many countries because it can influence people without them knowing.}
{Trong quảng cáo, đây là thông điệp tiềm thức.}

\vocabentry{Supplement}
{/ˈsʌplɪmənt/ (n)}
{A separate section of a newspaper or magazine.}
{The weekend edition of the paper includes a large travel supplement.}
{Đây là một danh từ dùng để chỉ tờ phụ trương.}

\vocabentry{Surveil}
{/sərˈveɪl/ (v)}
{Keep (a person or place) under surveillance.}
{Many people are concerned about how tech companies surveil their online activity.}
{Sự giám sát dữ liệu trong kỷ nguyên số.}

\vocabentry{Survey}
{/ˈsɜːrveɪ/ (n)}
{A general view, examination, or description of someone or something.}
{The survey results show that most people get their news from social media.}
{Đây là một danh từ dùng để chỉ cuộc khảo sát ý kiến độc giả/khách hàng.}

\vocabentry{Symmetric}
{/sɪˈmetrɪk/ (adj)}
{Made up of exactly similar parts facing each other or around an axis.}
{A symmetric layout can create a sense of balance and professionalism in an advertisement.}
{Trong thiết kế quảng cáo, đây là sự đối xứng.}

\vocabentry{Syndicate}
{/ˈsɪndɪkeɪt/ (v)}
{Publish or broadcast (material) simultaneously in a number of newspapers or television stations.}
{Her column is syndicated in over fifty different newspapers across the country.}
{Đây là một động từ dùng để chỉ việc bán nội dung cho nhiều đài/báo cùng lúc.}

\vocabentry{Talk show}
{/ˈtɔːk ʃoʊ/ (n)}
{A television or radio program in which celebrities and other guests are interviewed.}
{Late-night talk shows often feature interviews with actors, musicians, and politicians.}
{Đây là một danh từ dùng để chỉ chương trình trò chuyện.}

\vocabentry{Teaser}
{/ˈtiːzər/ (n)}
{A short introductory advertisement for a product, especially one that does not mention the name of the product.}
{The movie studio released a thirty-second teaser to build hype for the new blockbuster.}
{Đây là một danh từ dùng để chỉ đoạn quảng cáo ngắn khơi gợi sự tò mò.}

\vocabentry{Telecommunications}
{/ˌtelikəˌmjuːnɪˈkeɪʃnz/ (n)}
{Communication over a distance by cable, telegraph, telephone, or broadcasting.}
{The telecommunications industry has revolutionized the way we share information.}
{Đây là một danh từ dùng để chỉ ngành viễn thông, bao gồm các hệ thống truyền thông tin từ xa}

\vocabentry{Theme}
{/θiːm/ (n)}
{An idea that recurs in or pervades a work of art or literature.}
{The common theme of the series is the struggle for survival in a post-apocalyptic world.}
{Trong truyền thông, đây là chủ đề cốt lõi.}

\vocabentry{Thought-provoking}
{/ˈθɔːt prəˌvoʊkɪŋ/ (adj)}
{Stimulating careful consideration or attention.}
{The documentary was very thought-provoking and sparked a lot of debate online.}
{Đây là một tính từ dùng để chỉ sự gợi suy nghĩ, trăn trở.}

\vocabentry{Time slot}
{/ˈtaɪm slɑːt/ (n)}
{A specific time assigned for a broadcast or advertisement.}
{The new comedy show has been given a prime-time time slot on Friday nights.}
{Đây là một cụm danh từ dùng để chỉ khung giờ phát sóng.}

\vocabentry{Topic}
{/ˈtɑːpɪk/ (n)}
{A matter dealt with in a text, discourse, or conversation; a subject.}
{The journalist chose a controversial topic to attract more readers to his column.}
{Đây là một danh từ dùng để chỉ chủ đề bài viết/tin tức.}

\vocabentry{Trace}
{/treɪs/ (v)}
{Find or discover by investigation.}
{The reporter managed to trace the anonymous tip back to a former employee of the company.}
{Trong báo chí, đây là việc truy vết nguồn gốc thông tin.}

\vocabentry{Transparency}
{/trænsˈpærənsi/ (n)}
{The condition of being transparent.}
{Media organizations should have more transparency about their funding and ownership.}
{Sự minh bạch trong báo chí và quảng cáo.}

\vocabentry{Trend}
{/trend/ (n)}
{A general direction in which something is developing or changing.}
{One of the latest trends in marketing is the use of short-form video content.}
{Đây là một danh từ dùng để chỉ xu hướng truyền thông/mạng xã hội.}

\vocabentry{Tune in}
{/ˌtuːn ˈɪn/ (v)}
{Adjust a radio or television to receive a signal from a particular station.}
{Millions of people tuned in to watch the final of the world cup.}
{Đây là một cụm động từ dùng để chỉ việc đón nghe/xem một chương trình.}

\vocabentry{Turn on}
{/ˌtɜːrn ˈɑːn/ (v)}
{Cause (an electrical device) to function by operating a switch.}
{I usually turn on the radio as soon as I get into the car.}
{Hành động cơ bản để tiếp cận truyền thông truyền thống.}

\vocabentry{Ultimate}
{/ˈʌltɪmət/ (adj)}
{Being or happening at the end of a process; final.}
{The ultimate goal of any advertising campaign is to increase sales and profits.}
{Mục tiêu cuối cùng của quảng cáo là doanh số.}

\vocabentry{Unbiased}
{/ʌnˈbaɪəst/ (adj)}
{Showing no prejudice for or against something; impartial.}
{It is difficult to find completely unbiased news in today's polarized political climate.}
{Đây là một tính từ dùng để chỉ tính không thiên vị, công tâm.}

\vocabentry{Undermine}
{/ˌʌndərˈmaɪn/ (v)}
{Lessen the effectiveness, power, or ability of, especially gradually or insidiously.}
{The spread of misinformation can undermine public trust in scientific experts.}
{Tin giả có thể làm suy yếu sự tin tưởng của công chúng.}

\vocabentry{Update}
{/ˌʌpˈdeɪt/ (v)}
{Make (something) more modern or up to date.}
{The news website is updated every few minutes with the latest headlines.}
{Đây là một động từ dùng để chỉ việc cập nhật tin tức.}

\vocabentry{Upload}
{/ˌʌpˈloʊd/ (v)}
{Transfer (data) from one computer system to another, typically to a larger computer or a server.}
{It only takes a few seconds to upload a photo to social media.}
{Đây là một động từ dùng để chỉ việc tải nội dung lên mạng.}

\vocabentry{Vibrant}
{/ˈvaɪbrənt/ (adj)}
{Full of energy and life.}
{The city has a vibrant independent media scene with many small magazines and websites.}
{Đây là một tính từ dùng để chỉ môi trường truyền thông sôi động.}

\vocabentry{Video}
{/ˈvɪdioʊ/ (n)}
{The recording, reproducing, or broadcasting of moving visual images.}
{Online video consumption is growing at an incredible rate.}
{Đây là một danh từ dùng để chỉ video, phương tiện truyền thông chính hiện nay.}

\vocabentry{Viewpoint}
{/ˈvjuːpɔɪnt/ (n)}
{A particular attitude or way of considering a matter.}
{The documentary presents a variety of different viewpoints on the issue of climate change.}
{Đây là một danh từ dùng để chỉ góc nhìn, quan điểm.}

\vocabentry{Virtual}
{/ˈvɜːrtʃuəl/ (adj)}
{Almost or nearly as described, but not completely.}
{Many news organizations now have a major virtual presence through their websites and apps.}
{Đây là một tính từ dùng để chỉ sự ảo/trực tuyến.}

\vocabentry{Visual}
{/ˈvɪʒuəl/ (adj)}
{Relating to seeing or sight.}
{Strong visual elements are essential for creating an eye-catching advertisement.}
{Yếu tố thị giác trong quảng cáo.}

\vocabentry{Vlog}
{/vlɑːɡ/ (n)}
{A personal website or social media account where a person regularly posts short videos.}
{She has a popular vlog where she shares her daily life and travels.}
{Đây là một danh từ dùng để chỉ nhật ký bằng video.}

\vocabentry{Voice}
{/vɔɪs/ (n)}
{An expressed opinion or desire.}
{The media should give a voice to marginalized groups in society.}
{Trong truyền thông, đây là tiếng nói/quan điểm của một nhóm người.}

\vocabentry{Volume}
{/ˈvɑːljuːm/ (n)}
{The amount of space that a substance or object occupies; the quantity or power of sound.}
{The sheer volume of news available online can be overwhelming.}
{Lượng thông tin hoặc âm lượng phát sóng.}

\vocabentry{Voucher}
{/ˈvaʊtʃər/ (n)}
{A small printed piece of paper that entitles the holder to a discount.}
{You can find discount vouchers for several local restaurants in the newspaper.}
{Công cụ khuyến mãi trong quảng cáo.}

\vocabentry{Walkthrough}
{/ˈwɔːkθruː/ (n)}
{A demonstration or explanation of something.}
{The website features a step-by-step walkthrough of how to use the new software.}
{Trong truyền thông, đây là bài hướng dẫn chi tiết.}

\vocabentry{Warp}
{/wɔːrp/ (v)}
{Become or cause to become bent or twisted out of shape.}
{Sensationalist reporting can warp our perception of reality.}
{Nghĩa bóng: làm sai lệch thông tin.}

\vocabentry{Watch}
{/wɑːtʃ/ (v)}
{Look at or observe attentively.}
{I like to watch the news while I'm eating breakfast.}
{Hành động tiêu thụ truyền hình/video.}

\vocabentry{Website}
{/ˈwebsaɪt/ (n)}
{A set of related web pages located under a single domain name.}
{Most newspapers now have a website where you can read their stories for free.}
{Đây là một danh từ dùng để chỉ trang web.}

\vocabentry{Weekly}
{/ˈwiːkli/ (adj)}
{Occurring or published once a week.}
{The Sunday Times is a famous weekly broadsheet in the UK.}
{Đây là một tính từ dùng để chỉ các ấn phẩm hàng tuần.}

\vocabentry{Whistleblowing}
{/ˈwɪslbloʊɪŋ/ (n)}
{The act of informing on a person or organization engaged in an illicit activity.}
{Whistleblowing is often the only way to expose deep-seated corruption in large organizations.}
{Đây là một danh từ dùng để chỉ hành động tố giác, phanh phui hoạt động bất hợp pháp của tổ chức hoặc cá nhân}

\vocabentry{Wide}
{/waɪd/ (adj)}
{Over a large area.}
{The story has received wide coverage in both national and international media.}
{Độ phủ rộng của truyền thông.}

\vocabentry{Widespread}
{/ˌwaɪdˈspred/ (adj)}
{Found or distributed over a large area or number of people.}
{There is widespread concern about the impact of social media on mental health.}
{Đây là một tính từ dùng để chỉ sự phổ biến rộng rãi, mang tính lan tỏa trên diện rộng}

\vocabentry{Window}
{/ˈwɪndoʊ/ (n)}
{A rectangular area on a computer screen that displays an app or a document.}
{You can open multiple windows to compare news from different sources at the same time.}
{Cửa sổ trình duyệt/ứng dụng truyền thông.}

\vocabentry{Witness}
{/ˈwɪtnəs/ (n)}
{A person who sees an event, typically a crime or accident, take place.}
{The reporter interviewed several witnesses to the accident.}
{Nhân chứng là nguồn tin quan trọng cho báo chí.}

\vocabentry{Witty}
{/ˈwɪti/ (adj)}
{Showing or characterized by quick and inventive verbal humor.}
{The advertisement features a witty script that makes the audience laugh.}
{Lối viết hóm hỉnh trong quảng cáo/bình luận.}

\vocabentry{Wordsmith}
{/ˈwɜːrdsmɪθ/ (n)}
{A skilled user of words.}
{Copywriters are professional wordsmiths who know how to craft a compelling message.}
{Đây là một danh từ dùng để chỉ người viết lách giỏi/bút chiến.}

